% Worked solutions for Listas/lista-r1-accounting-capital-structure.md (15 exercises).
% Compile: pdflatex twice; check: pdffonts (Type 1 only, uni = yes).
% Every number below is asserted in lista-r1-accounting-capital-structure-solutions_check.py.
\documentclass[11pt,a4paper]{article}
\usepackage[utf8]{inputenc}
\usepackage[T1]{fontenc}
\usepackage{lmodern}   % Type 1 vector fonts (replaces PK bitmaps)
\usepackage{cmap}      % ToUnicode CMap -> copyable text and accents
\usepackage[brazilian,english]{babel}
\usepackage{amsmath,amssymb,amsfonts}
\usepackage[a4paper,margin=2.2cm]{geometry}
\usepackage{booktabs}
\usepackage{enumitem}
\usepackage{fancyhdr}
\usepackage{xcolor}
\usepackage{microtype}

\pagestyle{fancy}\fancyhf{}
\lhead{\footnotesize Finan\c{c}as II --- Corporativas (MPE Insper)}
\rhead{\footnotesize List R1 --- Worked solutions}
\cfoot{\footnotesize\thepage}
\renewcommand{\headrulewidth}{0.3pt}
\setlength{\headheight}{14pt}
\setlist{topsep=2pt,itemsep=1pt}

\newcommand{\exsec}[1]{\section*{\color{blue!55!black}#1}\addcontentsline{toc}{section}{#1}}
\newcommand{\notation}{\par\noindent\textbf{Notation.}}
\newcommand{\step}[2]{\par\medskip\noindent\textbf{Step #1 --- #2.}~}
\newcommand{\intuition}{\par\medskip\noindent\textbf{Intuition.}~}
\newcommand{\bookans}{\par\medskip\noindent\textbf{Book answer.}~}
\newcommand{\refbib}{\par\smallskip\noindent\textbf{Ref:}~}
\newcommand{\assume}{\par\smallskip\noindent\textbf{Assumption.}~}
\newcommand{\pt}{\foreignlanguage{brazilian}}

\title{List R1 --- Introduction, accounting revision and the capital structure\\[2pt]
\large Worked solutions}
\author{Finan\c{c}as II --- Corporativas (MPE Insper)}
\date{}

\begin{document}
\maketitle
\noindent Statements are in \texttt{Listas/lista-r1-accounting-capital-structure.md}; each
exercise is cited by number and printed page and is not restated. Editions are the local
copies: Higgins = Koski, Higgins \& Mitton, \emph{Analysis for Financial Management}, 13th ed.
(2023); Berk = Berk, DeMarzo \& Harford, \emph{Fundamentals of Corporate Finance}, 2nd ed.
(2011); Ross = Ross, Westerfield, Jordan \& Lamb, \emph{Fundamentos de Administra\c{c}\~ao
Financeira}, 9th ed. (2013). Ross prints thousands with a period (\pt{\$196.000}); below they are
written $196{,}000$. Money is in the units of each statement.

\tableofcontents

% =====================================================================
\exsec{1. Transactions and the balance sheet --- Higgins 1.4, p.~33}

\notation
\begin{itemize}
  \item $A$, $L$, $E$: total assets, total liabilities, owners' equity (\$; $A = L + E$).
  \item $\Delta X$: change in account $X$ caused by the transaction (\$, any sign).
  \item $C_f$, $C_g$: book (inventory) cost of the merchandise sold in (f) and (g) (\$, $\ge 0$; not given).
\end{itemize}
\assume Taxes are ignored (a gain or profit goes to retained earnings in full); in (e) the
repurchased shares are carried as treasury stock, a deduction from equity.

\step{1}{The identity in changes}
The balance sheet identity holds before and after each transaction,
\begin{align}
  A &= L + E, \label{e1-id}\\
  \Delta A &= \Delta L + \Delta E. \label{e1-did}
\end{align}

\step{2}{(a) Sale of used equipment}
Cash rises by the price and equipment falls by its book value:
\begin{align}
  \Delta A &= \Delta\text{Cash} + \Delta\text{Equipment} \notag\\
  &= (500{,}000) + (-300{,}000) \notag\\
  &= 200{,}000. \label{e1-a}
\end{align}
By \eqref{e1-did}, $\Delta E = 200{,}000$: the gain on sale (price minus book value) enters
retained earnings; liabilities do not change.

\step{3}{(b) Building financed 40\% cash, 60\% loan}
\begin{align}
  \Delta\text{Cash} &= -(0.40)(80) \notag\\
  &= -32 \ \text{million}, \label{e1-b1}\\
  \Delta\text{Bank loan} &= (0.60)(80) \notag\\
  &= 48 \ \text{million}. \label{e1-b2}
\end{align}
Assets change by $\Delta A = (80) + (-32) = 48$ million, equal to $\Delta L = 48$ from
\eqref{e1-b2}; equity is unchanged.

\step{4}{(c) Building bought for cash}
$\Delta A = (60) + (-60) = 0$: buildings $+60$ million, cash $-60$ million. Totals are
unchanged; only the composition of assets moves (less liquid).

\step{5}{(d) Payment to trade creditors}
Cash $-40{,}000$ and accounts payable $-40{,}000$: $\Delta A = \Delta L = -40{,}000$, $\Delta E = 0$.

\step{6}{(e) Share repurchase}
\begin{align}
  \Delta\text{Cash} &= -(10{,}000)(24) \notag\\
  &= -240{,}000. \label{e1-e}
\end{align}
Equity falls by the same amount (treasury stock), $\Delta E = -240{,}000$; liabilities unchanged.

\step{7}{(f) Cash sale of merchandise}
Cash $+80{,}000$, inventory $-C_f$, so $\Delta A = 80{,}000 - C_f$ and, by \eqref{e1-did}, retained
earnings rise by the gross profit $\Delta E = 80{,}000 - C_f$.

\step{8}{(g) Credit sale of merchandise}
Accounts receivable $+120{,}000$, inventory $-C_g$, $\Delta E = 120{,}000 - C_g$. Same effect on
equity as a cash sale; the asset received is a receivable, not cash.

\step{9}{(h) Dividend}
Cash $-50{,}000$ and retained earnings $-50{,}000$: $\Delta A = \Delta E = -50{,}000$.

\intuition Every transaction has two sides. Only exchanges with outsiders that create or destroy
value for owners (a gain, a profit, a payout, a buyback) move equity; swaps of one asset for
another, or borrowing to buy an asset, change the composition or size of the balance sheet
without touching equity.

\refbib Higgins Ch.~1, \S\,``The Balance Sheet'' (with ``Current Assets and Liabilities''),
pp.~9--11; \S\,``Measuring Earnings'', p.~12. No printed answer (even-numbered problem).

% =====================================================================
\exsec{2. Local Co.\ margins --- Berk 2.12, p.~54}

\notation
\begin{itemize}
  \item $S$: sales; $C$: cost of sales; $SGA$: selling, general and administrative expenses;
        $RD$: research and development; $D$: depreciation (\$ million, $\ge 0$).
  \item $\tau$: tax rate ($0 \le \tau < 1$); $GP$: gross profit; $OI$: operating income; $NI$: net income.
\end{itemize}
\assume No interest expense and no other income (none is given), so $OI$ is also pretax income.
Depreciation is an operating expense, as in Berk's Table 2.3 (p.~34).

\step{1}{Gross profit and gross margin (a)}
\begin{align}
  GP &= S - C \label{e2-gp}\\
  &= (10) - (6) \notag\\
  &= 4, \notag\\
  \text{Gross margin} &= GP/S \label{e2-gm}\\
  &= (4)/(10) \notag\\
  &= 0.40 \notag\\
 &= 40\%. \notag
\end{align}

\step{2}{Operating income and operating margin (b)}
Using \eqref{e2-gp},
\begin{align}
  OI &= GP - SGA - RD - D \label{e2-oi}\\
  &= (4) - (0.5) - (1) - (1) \notag\\
  &= 1.5, \notag\\
  \text{Operating margin} &= OI/S \notag\\
  &= (1.5)/(10) \notag\\
  &= 0.15 \notag\\
 &= 15\%. \notag
\end{align}

\step{3}{Net income and net profit margin (c)}
Using \eqref{e2-oi} and no interest,
\begin{align}
  NI &= OI(1-\tau) \label{e2-ni}\\
  &= (1.5)(1-0.35) \notag\\
  &= 0.975, \notag\\
  \text{Net profit margin} &= NI/S \notag\\
  &= (0.975)/(10) \notag\\
  &= 0.0975 \notag\\
 &= 9.75\%. \notag
\end{align}

\intuition Each margin peels off one more layer of cost: production costs (gross), the cost of
running the business (operating), then financing and taxes (net). The gap between operating and
net margin here is only the 35\% tax.

\refbib Berk Ch.~2, \S2.4 ``The Income Statement'', pp.~33--35; \S2.5 ``Income Statement Analysis''
(Eqs.~2.9--2.11), pp.~35--36.

% =====================================================================
\exsec{3. Nokela: earnings versus cash flow --- Berk 2.31, p.~57}

\notation
\begin{itemize}
  \item $I = 40$: purchase price (\$ million), paid in year 1 (``this year'').
  \item $D = 10$: straight-line depreciation in each year $t = 1,\dots,4$ (\$ million).
  \item $\tau = 0.40$: tax rate; $\Delta NI_t$, $\Delta CF_t$: change in earnings and in cash flow in year $t$.
\end{itemize}
\assume Nokela has enough other taxable income to use the depreciation deduction each year.

\step{1}{Impact on earnings (a)}
Depreciation is deducted before tax, so each year it lowers net income by its after-tax amount:
\begin{align}
  \Delta NI_t &= -D(1-\tau) \label{e3-ni}\\
  &= -(10)(1-0.40) \notag\\
  &= -6 \quad (t \notag\\
 &= 1,2,3,4). \notag
\end{align}

\step{2}{Impact on cash flow (b)}
The purchase is a cash outflow in year 1; depreciation is not a cash cost, but it saves taxes of
$\tau D$ each year:
\begin{align}
  \tau D &= (0.40)(10) \notag\\
  &= 4, \label{e3-ts}\\
  \Delta CF_1 &= -I + \tau D \notag\\
  &= -(40) + (4) \notag\\
  &= -36, \label{e3-cf1}\\
  \Delta CF_t &= \tau D = 4 \quad (t \notag\\
  &= 2,3,4). \label{e3-cft}
\end{align}

\step{3}{Consistency check}
Over the four years both measures add to the same total, using \eqref{e3-ni}--\eqref{e3-cft}:
\begin{align}
  \textstyle\sum_t \Delta NI_t &= 4(-6) \notag\\
  &= -24, \notag\\
  \textstyle\sum_t \Delta CF_t &= (-36) + 3(4) \notag\\
  &= -24. \notag
\end{align}

\intuition Accounting spreads the cost evenly ($-6$ a year); cash flow records it when paid
($-36$ up front) and then gives back the tax shield ($+4$ a year). Same total, different timing,
and timing is what matters for value.

\refbib Berk Ch.~2, \S2.4, pp.~33--35 (depreciation as a non-cash expense); \S2.6 ``The Statement
of Cash Flows'', pp.~42--45.

% =====================================================================
\exsec{4. GE: book versus market --- Berk 2.8, p.~53}

\notation
\begin{itemize}
  \item $N = 10.3$ billion shares; $P = 38$: share price (\$); $BE = 117$: book equity (\$ billion).
  \item $D = 467$: total debt; $Cash = 16$ (\$ billion); $MV_E$: market capitalization; $EV$: enterprise value.
\end{itemize}

\step{1}{Market capitalization and market-to-book (a)}
\begin{align}
  MV_E &= N \times P \label{e4-mc}\\
  &= (10.3)(38) \notag\\
  &= 391.4, \notag\\
  M/B &= MV_E/BE \notag\\
  &= (391.4)/(117) \notag\\
  &= 3.345. \notag
\end{align}

\step{2}{Debt-equity ratios (b)}
\begin{align}
  (D/E)_{\text{book}} &= D/BE \notag\\
  &= (467)/(117) \notag\\
  &= 3.991, \notag\\
  (D/E)_{\text{market}} &= D/MV_E \notag\\
  &= (467)/(391.4) \notag\\
  &= 1.193. \notag
\end{align}

\step{3}{Enterprise value (c)}
Using \eqref{e4-mc} and Berk's Eq.~2.5,
\begin{align}
  EV &= MV_E + D - Cash \notag\\
  &= (391.4) + (467) - (16) \notag\\
  &= 842.4 \ \text{billion}. \notag
\end{align}

\intuition Investors valued GE's equity at more than three times its book value, so leverage
measured at book (4.0) overstates the claim of debt relative to what equity is worth (1.2).
Enterprise value nets out cash: it is the value of the operating business that a buyer of all the
equity and all the debt would acquire.

\refbib Berk Ch.~2, \S2.3 ``Balance Sheet Analysis'' (market-to-book, debt-equity, Eq.~2.5),
pp.~30--33.

% =====================================================================
\exsec{5. True or false on ratios --- Higgins 2.3, p.~71}

\notation
\begin{itemize}
  \item $A$, $L$, $E$: assets, liabilities, equity (\$; $L \ge 0$, $E > 0$ unless stated).
  \item $NI$: net income (any sign); $CA$, $CL$, $Inv$: current assets, current liabilities, inventory.
  \item $P$: share price; $EPS$: earnings per share.
\end{itemize}

\step{1}{(a) True}
\begin{align}
  \frac{A}{E} &= \frac{L + E}{E} \notag\\
  &= 1 + \frac{L}{E}. \label{e5-a}
\end{align}
This is the balance sheet identity divided by $E$.

\step{2}{(b) False as stated; true for non-negative income}
Divide $ROE = NI/E$ by $ROA = NI/A$:
\begin{align}
  ROE &= ROA \times \frac{A}{E}. \label{e5-b}
\end{align}
By \eqref{e5-a}, $A/E \ge 1$ when $L \ge 0$. If $NI \ge 0$, \eqref{e5-b} gives $ROE \ge ROA$. If
$NI < 0$, the same multiplier makes $ROE$ \emph{more} negative: with $NI = -10$, $A = 100$,
$E = 50$, $ROA = -10\%$ and $ROE = (-10\%)(2) = -20\% < ROA$. Leverage amplifies returns in both
directions, so ``always'' fails for loss-making firms (and the comparison breaks with negative equity).

\step{3}{(c) False}
The collection period and the payables period are set largely by industry practice and bargaining
power. A payables period longer than the collection period is convenient (suppliers finance the
receivables) but nothing forces it.

\step{4}{(d) True}
\begin{align}
  \frac{CA}{CL} - \frac{CA - Inv}{CL} &= \frac{Inv}{CL} \ge 0, \notag
\end{align}
since inventory is never negative and $CL > 0$.

\step{5}{(e) True}
By the DuPont decomposition $ROE = (NI/S)(S/A)(A/E)$, a higher asset turnover $S/A$ with margin
and leverage held fixed raises ROE: the same sales are produced with fewer assets.

\step{6}{(f) False}
The earnings yield is $EPS/P = 1/(P/EPS)$; equal earnings yields imply equal P/E ratios.

\step{7}{(g) False}
Without taxes and transaction costs, a paper gain can be turned into cash at any moment by
selling, so it is worth exactly as much as an equal realized gain.

\intuition These are accounting identities, not empirical regularities: (a), (d), (f) follow
from definitions; (b) holds only because leverage is a multiplier, which also magnifies losses.

\bookans Matches the book's answers (p.~401) for (a), (c)--(g). For (b) the book answers True,
arguing that ROA can exceed ROE only with negative liabilities; that argument implicitly assumes
$NI \ge 0$. With a net loss, $ROE < ROA$ by \eqref{e5-b}, so I mark (b) False as stated.

\refbib Higgins Ch.~2, \S\,``The Three Determinants of ROE'', pp.~38--40; ``Asset Turnover'',
pp.~42--48; ``Financial Leverage'', pp.~48--53; ``The Value Problem'', pp.~56--58; Ch.~1,
``Economic Income vs.\ Accounting Income'', pp.~26--27.

% =====================================================================
\exsec{6. Book and market leverage of two firms --- Berk 2.21, p.~55}

\notation
\begin{itemize}
  \item $D$: debt; $BE$: book equity; $ME$: market equity; $OI$: operating income; $Int$: interest
        expense (\$ million, $> 0$); subscripts $A$, $B$ for the firm.
\end{itemize}
\assume Interest coverage uses operating income, one of the earnings measures Berk lists (p.~37).

\step{1}{Market debt-equity (a)}
\begin{align}
  (D/E)^{mkt}_A &= (500)/(400) \notag\\
  &= 1.25, \notag\\
  (D/E)^{mkt}_B &= (80)/(40) \notag\\
  &= 2.00. \notag
\end{align}

\step{2}{Book debt-equity (b)}
\begin{align}
  (D/E)^{bk}_A &= (500)/(300) \notag\\
  &= 1.667, \notag\\
  (D/E)^{bk}_B &= (80)/(35) \notag\\
  &= 2.286. \notag
\end{align}

\step{3}{Interest coverage (c)}
\begin{align}
  \text{Coverage} &= OI/Int, \label{e6-cov}\\
  \text{Coverage}_A &= (100)/(50) \notag\\
  &= 2.00, \notag\\
  \text{Coverage}_B &= (8)/(7) \notag\\
  &= 1.143. \notag
\end{align}

\step{4}{Who struggles more (d)}
Firm B: higher leverage on both measures, and by \eqref{e6-cov} its operating income covers interest
only 1.14 times, so a fall of about 12.5\% in operating income ($1/8$ of it) would leave it unable to pay
interest out of earnings; Firm A can absorb a 50\% fall.

\intuition Ratios scale away firm size. Firm B is about a sixth of A's size but carries
proportionally more debt and has almost no cushion of earnings above its interest bill.

\refbib Berk Ch.~2, \S2.3 (debt-equity ratios), pp.~30--33; \S2.5 ``Leverage Ratios'', p.~37.

% =====================================================================
\exsec{7. DuPont for a retail firm --- Berk 2.27, p.~56}

\notation
\begin{itemize}
  \item $m$: net profit margin $NI/S$; $T$: total asset turnover $S/A$; $A = 44$: total assets;
        $BE = 18$: book equity (\$ million); $EM = A/BE$: equity multiplier.
\end{itemize}
\assume In (c) book equity stays at 18 (the statement fixes assets and liabilities), so only $T$ changes.

\step{1}{Equity multiplier and DuPont identity}
\begin{align}
  ROE &= m \times T \times EM, \label{e7-dp}\\
  EM &= (44)/(18) \notag\\
  &= 2.4444. \label{e7-em}
\end{align}

\step{2}{Current ROE (a)}
Substituting \eqref{e7-em} into \eqref{e7-dp},
\begin{align}
  ROE_a &= (0.035)(1.8)(2.4444) \notag\\
  &= 0.1540 \notag\\
 &= 15.4\%. \notag
\end{align}
Check in levels: $S = (1.8)(44) = 79.2$, $NI = (0.035)(79.2) = 2.772$, $ROE = (2.772)/(18) = 0.154$.

\step{3}{Higher margin (b)}
\begin{align}
  ROE_b &= (0.04)(1.8)(2.4444) \notag\\
  &= 0.1760 \notag\\
 &= 17.6\%. \notag
\end{align}

\step{4}{Higher margin and 20\% more revenue (c)}
With $A$ fixed, turnover rises by 20\%:
\begin{align}
  T_c &= (1.8)(1.20) \notag\\
  &= 2.16, \label{e7-tc}\\
  ROE_c &= (0.04)(2.16)(2.4444) \notag\\
  &= 0.2112 \notag\\
 &= 21.1\%. \notag
\end{align}
In levels: $S = 95.04$, $NI = (0.04)(95.04) = 3.8016$, $ROE = 3.8016/18 = 0.2112$.

\intuition ROE is a product, so proportional changes multiply: $+14.3\%$ in margin
($4/3.5$) and $+20\%$ in turnover give $(1.143)(1.2) = 1.371$ times the original ROE.

\refbib Berk Ch.~2, \S2.5 ``Investment Returns'' and ``The DuPont Identity'' (Eq.~2.18), pp.~37--39.

% =====================================================================
\exsec{8. Atlantic vs.\ Pacific: ROE, ROA, ROIC --- Higgins 2.7, p.~73}

\notation
\begin{itemize}
  \item $EBIT$: earnings before interest and taxes; $D$: debt; $E$: equity (\$ million); $r = 0.07$:
        interest rate; $t = 0.30$: tax rate.
  \item $A = D + E$: total assets (no other liabilities); $NI$: net income.
  \item $ROE = NI/E$; $ROA = NI/A$; $ROIC = EBIT(1-t)/(D+E)$ (Higgins' definitions).
\end{itemize}

\step{1}{Net income}
\begin{align}
  NI &= (EBIT - rD)(1-t), \label{e8-ni}\\
  NI_{Atl} &= (450 - (0.07)(290))(0.70) \notag\\
  &= (450 - 20.3)(0.70) \notag\\
  &= (429.7)(0.70) \notag\\
  &= 300.79, \notag\\
  NI_{Pac} &= (470 - (0.07)(1490))(0.70) \notag\\
  &= (470 - 104.3)(0.70) \notag\\
  &= (365.7)(0.70) \notag\\
  &= 255.99. \notag
\end{align}

\step{2}{ROE, ROA, ROIC (a)}
Using \eqref{e8-ni} and $A_{Atl} = 1200$, $A_{Pac} = 1860$:
\begin{align}
  ROE_{Atl} &= (300.79)/(910) \notag\\
  &= 0.3305, \notag\\
  ROA_{Atl} &= (300.79)/(1200) \notag\\
  &= 0.2507, \notag\\
  ROIC_{Atl} &= (450)(0.70)/(1200) \notag\\
  &= 0.2625, \notag\\
  ROE_{Pac} &= (255.99)/(370) \notag\\
  &= 0.6919, \notag\\
  ROA_{Pac} &= (255.99)/(1860) \notag\\
  &= 0.1376, \notag\\
  ROIC_{Pac} &= (470)(0.70)/(1860) \notag\\
  &= 0.1769. \notag
\end{align}
So ROE 33.1\% vs.\ 69.2\%, ROA 25.1\% vs.\ 13.8\%, ROIC 26.3\% vs.\ 17.7\%.

\step{3}{Why Pacific's ROE is higher (b)}
Write $NI = [ROIC(D+E) - r(1-t)D]$, divide by $E$:
\begin{align}
  ROE &= ROIC + \bigl(ROIC - r(1-t)\bigr)\frac{D}{E}. \label{e8-lev}
\end{align}
With $r(1-t) = (0.07)(0.70) = 0.049$, \eqref{e8-lev} gives
\begin{align}
  ROE_{Atl} &= (0.2625) + (0.2625 - 0.049)(290/910) \notag\\
  &= 0.3305, \notag\\
  ROE_{Pac} &= (0.1769) + (0.1769 - 0.049)(1490/370) \notag\\
  &= 0.6919, \notag
\end{align}
reproducing Step~2. Pacific's ROE comes from its debt-equity ratio of 4.03 (vs.\ 0.32), not from
better operations; it is also riskier for shareholders. It is not a better company on this evidence.

\step{4}{Why Atlantic's ROA is higher (c)}
ROA puts the income of shareholders only (after interest) over the capital of both shareholders and
creditors, so it penalizes leverage: Pacific's large interest bill lowers the numerator. The gap mixes
leverage with operating performance and says little on its own.

\step{5}{ROIC comparison (d)}
ROIC is independent of financing: Atlantic earns 26.3\% after tax on all invested capital against
Pacific's 17.7\%, so Atlantic's assets are more productive. Both exceed the after-tax cost of debt,
which is why leverage raises ROE in \eqref{e8-lev}; whether either beats its full cost of capital
depends on business risk, not shown here.

\intuition ROE answers ``how much per dollar of owners' money'', ROA mixes in the financing choice,
and ROIC isolates the business. Compare operating performance with ROIC, then read ROE as ROIC
plus a leverage bonus (or penalty).

\bookans Matches the book's answer (pp.~402--403): ROE 33.1\% / 69.2\%, ROA 25.1\% / 13.8\%,
ROIC 26.3\% / 17.7\%; same conclusions for (b)--(d).

\refbib Higgins Ch.~2, \S\,``The Three Determinants of ROE'', pp.~38--40; ``The Risk Problem''
(ROA, ROIC), pp.~54--56.

% =====================================================================
\exsec{9. Nave Espacial Jetson: cash flows --- Ross 2.14, p.~48}

\notation
\begin{itemize}
  \item $S$: sales; $C$: costs; $X$: other expenses; $Dep$: depreciation; $Int$: interest; $Tax$:
        taxes; $Div$: dividends (\$).
  \item $OCF$: operating cash flow (\pt{fluxo de caixa operacional}); $NCS$: net capital spending;
        $\Delta NWC$: change in net working capital (\pt{CCL}); $CFA$: cash flow from assets;
        $CFC$, $CFS$: cash flow to creditors and to shareholders.
\end{itemize}

\step{1}{EBIT and operating cash flow (a)}
\begin{align}
  EBIT &= S - C - X - Dep \notag\\
  &= (196{,}000) - (104{,}000) - (6{,}800) - (9{,}100) \notag\\
  &= 76{,}100, \label{e9-ebit}\\
  OCF &= EBIT + Dep - Tax \label{e9-ocfdef}\\
  &= (76{,}100) + (9{,}100) - (21{,}455) \notag\\
  &= 63{,}745. \label{e9-ocf}
\end{align}

\step{2}{Cash flow to creditors (b)}
Net new long-term borrowing was $-7{,}300$ (debt retired):
\begin{align}
  CFC &= Int - \text{net new borrowing} \label{e9-cfcdef}\\
  &= (14{,}800) - (-7{,}300) \notag\\
  &= 22{,}100. \label{e9-cfc}
\end{align}

\step{3}{Cash flow to shareholders (c)}
\begin{align}
  CFS &= Div - \text{net new equity} \notag\\
  &= (10{,}400) - (5{,}700) \notag\\
  &= 4{,}700. \label{e9-cfs}
\end{align}

\step{4}{Change in net working capital (d)}
The cash flow identity $CFA = CFC + CFS = OCF - NCS - \Delta NWC$ with
\eqref{e9-cfc}--\eqref{e9-cfs}:
\begin{align}
  CFA &= (22{,}100) + (4{,}700) \notag\\
  &= 26{,}800, \label{e9-cfa}\\
  NCS &= \Delta\text{NFA} + Dep \notag\\
  &= (27{,}000) + (9{,}100) \notag\\
  &= 36{,}100, \label{e9-ncs}\\
  \Delta NWC &= OCF - NCS - CFA \notag\\
  &= (63{,}745) - (36{,}100) - (26{,}800) \notag\\
  &= 845. \notag
\end{align}

\intuition The firm's operations threw off 63,745; it reinvested 36,100 in fixed assets and 845 in
working capital, and paid the remaining 26,800 to its investors, mostly to creditors (interest plus
debt repayment).

\refbib Ross Cap.~2, \S2.4 ``Fluxo de caixa'', pp.~37--44.

% =====================================================================
\exsec{10. Blake's Restaurant Supply: from statements to cash --- Higgins 1.11, p.~34}

\notation
\begin{itemize}
  \item Subscripts $0$ = end-2020, $1$ = end-2021; $S$: net sales 2021; $COGS$: cost of goods sold
        2021; $Dep$: depreciation 2021; $NI$: net income 2021 (\$ million).
  \item $AR$, $Inv$, $AP$, $NFA$, $Cash$: balance sheet accounts; $Capex$: capital expenditures;
        $CFO$: cash flow from operations.
\end{itemize}

\step{1}{Cash collected from sales (a)}
Receivables accumulate sales not yet collected: $AR_1 = AR_0 + S - \text{Collections}$, so
\begin{align}
  \text{Collections} &= S - (AR_1 - AR_0) \notag\\
  &= (782) - (87 - 57) \notag\\
  &= 752. \notag
\end{align}

\step{2}{Cost of goods produced (b)}
$Inv_1 = Inv_0 + \text{Production} - COGS$, so
\begin{align}
  \text{Production} &= COGS + (Inv_1 - Inv_0) \notag\\
  &= (502) + (29 - 39) \notag\\
  &= 492. \notag
\end{align}

\step{3}{Capital expenditures (c)}
$NFA_1 = NFA_0 + Capex - Dep$ (no disposals), so
\begin{align}
  Capex &= (NFA_1 - NFA_0) + Dep \notag\\
  &= (482 - 404) + (61) \notag\\
  &= 139. \label{e10-capex}
\end{align}

\step{4}{Cash flow from operations (d)}
With no financing flows, the change in cash is operating plus investing cash flow, and investing
is $-Capex$ from \eqref{e10-capex}:
\begin{align}
  \Delta Cash &= CFO - Capex, \notag\\
  CFO &= \Delta Cash + Capex \notag\\
  &= (135 - 86) + (139) \notag\\
  &= 188. \label{e10-cfo}
\end{align}
Indirect check from the income statement:
\begin{align}
  CFO &= NI + Dep - \Delta AR - \Delta Inv + \Delta AP \notag\\
  &= (142) + (61) - (30) - (-10) + (5) \notag\\
  &= 188. \notag
\end{align}

\intuition Accrual accounting records sales when made and costs when goods are sold; cash
records them when money moves. Rising receivables keep 30 of sales out of the bank; running down
inventory and stretching payables put cash back in.

\bookans Matches the book's answer (pp.~399--400): 752, 492, 139, 188 (\$ million).

\refbib Higgins Ch.~1, \S\,``Sources and Uses Statements'', pp.~16--18; ``The Cash Flow
Statement'', pp.~18--22.

% =====================================================================
\exsec{11. Ind\'ustrias Dahlia: cash flow identity --- Ross 2.21, p.~49}

\notation
\begin{itemize}
  \item $S$, $COGS$, $Dep$, $Int$, $Div$: sales, cost of goods sold, depreciation, interest,
        dividends (\$); $\tau = 0.34$.
  \item $NFA$, $CA$, $CL$: net fixed assets, current assets, current liabilities; subscripts $0$
        (start of year), $1$ (end of year). Other symbols as in Exercise 9.
\end{itemize}

\step{1}{Net income (a)}
\begin{align}
  EBIT &= (22{,}800) - (16{,}050) - (4{,}050) \notag\\
  &= 2{,}700, \label{e11-ebit}\\
  EBT &= (2{,}700) - (1{,}830) \notag\\
  &= 870, \notag\\
  Tax &= (0.34)(870) \notag\\
  &= 295.80, \label{e11-tax}\\
  NI &= (870) - (295.80) \notag\\
  &= 574.20. \notag
\end{align}

\step{2}{Operating cash flow (b)}
By \eqref{e9-ocfdef} with \eqref{e11-ebit} and \eqref{e11-tax},
\begin{align}
  OCF &= (2{,}700) + (4{,}050) - (295.80) \notag\\
  &= 6{,}454.20. \label{e11-ocf}
\end{align}

\step{3}{Cash flow from assets (c)}
\begin{align}
  NCS &= (16{,}800 - 13{,}650) + (4{,}050) \notag\\
  &= 7{,}200, \label{e11-ncs}\\
  \Delta NWC &= (5{,}930 - 3{,}150) - (4{,}800 - 2{,}700) \notag\\
  &= (2{,}780) - (2{,}100) \notag\\
  &= 680, \label{e11-nwc}\\
  CFA &= OCF - NCS - \Delta NWC \notag\\
  &= (6{,}454.20) - (7{,}200) - (680) \notag\\
  &= -1{,}425.80. \label{e11-cfa}
\end{align}
A negative CFA is possible: the firm invested $7{,}880$ in fixed assets and working capital, more
than the $6{,}454.20$ its operations generated, and raised the difference from investors.

\step{4}{Cash flow to creditors and shareholders (d)}
No new debt, so by \eqref{e9-cfcdef}
\begin{align}
  CFC &= (1{,}830) - (0) \notag\\
  &= 1{,}830, \label{e11-cfc}\\
  CFS &= CFA - CFC \notag\\
  &= (-1{,}425.80) - (1{,}830) \notag\\
  &= -3{,}255.80, \label{e11-cfs}\\
  \text{net new equity} &= Div - CFS \notag\\
  &= (1{,}300) - (-3{,}255.80) \notag\\
  &= 4{,}555.80. \notag
\end{align}

\intuition Signs: net income and OCF are positive (the business is profitable and generates cash);
CFA is negative (growth spending exceeds that cash); CFC is positive (creditors received their
interest); CFS is negative (shareholders put in 4,555.80 of new equity, net of the 1,300 dividend
they received). A growing firm is a net user of investors' cash.

\refbib Ross Cap.~2, \S2.2 ``A demonstra\c{c}\~ao de resultados do exerc\'icio'', pp.~30--34;
\S2.4 ``Fluxo de caixa'', pp.~37--44.

% =====================================================================
\exsec{12. Global's marketing campaign --- Berk *2.19, p.~54}

\notation
\begin{itemize}
  \item $S$: net sales; $m$: operating margin; $EBIT$; $Int$: interest expense; $\tau$: taxes /
        pretax income; $NI$: net income (\$ million); $N$: shares (million); $P$: share price (\$).
\end{itemize}
\textbf{Data used} (Berk Table 2.3, p.~34; Eq.~2.8, p.~34; p.~39). Reported 2010: net sales
$186.7$, operating income $=$ EBIT $= 10.4$ (no other income), so $m = 10.4/186.7 = 5.57\%$;
interest expense $7.7$; pretax income $2.7$; taxes $0.7$; net income $2.0$; $N = 3.6$ million
shares; P/E $= 36/2.0 = 18$ (p.~39). Reported 2009: pretax income $2.5$, taxes $0.6$.

\assume The 5.57\% margin is Global's reported 2010 margin, so the campaign scenario replaces the
reported 2010: sales are 15\% above the reported $186.7$ and interest stays at the reported $7.7$.
The tax rate is the 2009 ratio, as the statement says. (If the 2010 ratio $0.7/2.7$ were meant
instead, the results are given at the end.)

\step{1}{EBIT (a)}
\begin{align}
  S' &= (186.7)(1.15) \notag\\
  &= 214.705, \label{e12-s}\\
  EBIT' &= m' S' \notag\\
  &= (0.045)(214.705) \notag\\
  &= 9.6617. \label{e12-ebit}
\end{align}

\step{2}{Net income (b)}
\begin{align}
  \tau_{2009} &= (0.6)/(2.5) \notag\\
  &= 0.24, \label{e12-tau}\\
  EBT' &= (9.6617) - (7.7) \notag\\
  &= 1.9617, \notag\\
  NI' &= (1.9617)(1 - 0.24) \notag\\
  &= 1.4909. \label{e12-ni}
\end{align}

\step{3}{Share price (c)}
With P/E and $N$ unchanged,
\begin{align}
  P' &= (P/E) \times NI'/N \notag\\
  &= (18)(1.4909/3.6) \notag\\
  &= (18)(0.41414) \notag\\
  &= 7.45. \label{e12-p}
\end{align}
\textit{Alternative tax rate.} With $\tau = 0.7/2.7 = 0.25926$: $NI' = (1.9617)(0.74074) = 1.4531$ and
$P' = (18)(1.4531/3.6) = 7.27$.

\intuition Sales up 15\% but margin down by a fifth leaves EBIT below the reported $10.4$. Because
interest is fixed, the fall in EBIT ($-7.1\%$) becomes a fall of about $27\%$ in pretax income: that
is financial leverage at work, and with a constant P/E the price falls in proportion to earnings,
from \$10 to about \$7.45.

\refbib Berk Ch.~2, \S2.4, pp.~33--35 (Table 2.3); \S2.5 ``Profitability Ratios'' and ``Valuation
Ratios'' (Eqs.~2.10, 2.19), pp.~35--36, 39.

% =====================================================================
\exsec{13. Quisco: develop in house or acquire with stock --- Berk *2.23, p.~55}

\notation
\begin{itemize}
  \item $N = 6{,}500$ million shares; $P = 18$ (\$); $EPS_0 = 0.80$: EPS without the technology.
  \item $R = 500$: in-house R\&D cost; $V = 900$: value of shares paid for the target (\$ million);
        $\tau = 0.35$; $\Delta N$: new shares issued.
\end{itemize}
\assume New shares are issued at the current price of \$18.

\step{1}{Baseline net income}
\begin{align}
  NI_0 &= EPS_0 \times N \notag\\
  &= (0.80)(6{,}500) \notag\\
  &= 5{,}200. \label{e13-ni}
\end{align}

\step{2}{In-house development (a)}
The R\&D expense lowers net income by its after-tax amount:
\begin{align}
  \Delta NI &= -R(1-\tau) \notag\\
  &= -(500)(0.65) \notag\\
  &= -325, \label{e13-dni}\\
  \Delta EPS_a &= (-325)/(6{,}500) \notag\\
  &= -0.05, \notag\\
  EPS_a &= (0.80) - (0.05) \notag\\
  &= 0.75. \notag
\end{align}

\step{3}{Acquisition with stock (b)}
\begin{align}
  \Delta N &= V/P \notag\\
  &= (900)/(18) \notag\\
  &= 50, \label{e13-dn}\\
  EPS_b &= NI_0/(N + \Delta N) \notag\\
  &= (5{,}200)/(6{,}500 + 50) \notag\\
  &= 0.7939, \notag\\
  \Delta EPS_b &= (0.7939) - (0.80) \notag\\
  &= -0.0061. \notag
\end{align}

\step{4}{Which is cheaper (c)}
The acquisition has the smaller effect on reported EPS ($-0.6$ cents vs.\ $-5$ cents), but it is not
cheaper. Existing shareholders give up \$900 million of value (the 50 million shares go to the
target's owners), against an after-tax cost of \$325 million (\$500 million pretax) for in-house
development: building is cheaper by \$575 million after tax, assuming both routes deliver the
same technology.

\intuition EPS measures how accounting allocates a cost, not the cost itself. Expensing R\&D hits
earnings at once; paying with stock spreads the same economic transfer as dilution, which EPS
records only as a small rise in the share count.

\refbib Berk Ch.~2, \S2.4 (EPS, Eq.~2.8), pp.~33--35; \S2.8 ``Financial Reporting in Practice'',
pp.~46--48.

% =====================================================================
\exsec{14. Complete the balance sheet --- Higgins 2.11, p.~74}

\notation
\begin{itemize}
  \item $S$: annual sales (all on credit); $COGS$: cost of goods sold (\$); 365-day year.
  \item $Cash$, $AR$, $Inv$, $CA$, $NFA$, $TA$: cash, receivables, inventory, current assets,
        net fixed assets, total assets; $AP$, $STD$, $CL$, $LTD$, $TL$, $E$: payables, short-term
        debt, current liabilities, long-term debt, total liabilities, equity.
  \item Ratios: days' sales in cash $= Cash/(S/365)$; collection period $= AR/(S/365)$; inventory
        turnover $= COGS/Inv$; payables period $= AP/(COGS/365)$; current ratio $= CA/CL$;
        liabilities to assets $= TL/TA$.
\end{itemize}

\step{1}{Sales from days' sales in cash}
\begin{align}
  S &= 365 \times Cash/34 \notag\\
  &= 365 \times (1{,}100{,}000)/34 \notag\\
  &= 11{,}808{,}824. \label{e14-s}
\end{align}

\step{2}{Receivables and current assets}
Using \eqref{e14-s},
\begin{align}
  AR &= 71 \times S/365 \notag\\
  &= 71 \times (11{,}808{,}824)/365 \notag\\
  &= 2{,}297{,}059, \label{e14-ar}\\
  CA &= (1{,}100{,}000) + (2{,}297{,}059) + (1{,}900{,}000) \notag\\
  &= 5{,}297{,}059, \label{e14-ca}\\
  NFA &= (8{,}000{,}000) - (5{,}297{,}059) \notag\\
  &= 2{,}702{,}941. \notag
\end{align}

\step{3}{COGS and payables}
\begin{align}
  COGS &= 5 \times Inv \notag\\
  &= 5 \times (1{,}900{,}000) \notag\\
  &= 9{,}500{,}000, \label{e14-cogs}\\
  AP &= 36 \times COGS/365 \notag\\
  &= 36 \times (9{,}500{,}000)/365 \notag\\
  &= 936{,}986. \label{e14-ap}
\end{align}

\step{4}{Current liabilities and short-term debt}
Using \eqref{e14-ca} and \eqref{e14-ap},
\begin{align}
  CL &= CA/2.6 \notag\\
  &= (5{,}297{,}059)/2.6 \notag\\
  &= 2{,}037{,}330, \label{e14-cl}\\
  STD &= CL - AP \notag\\
  &= (2{,}037{,}330) - (936{,}986) \notag\\
  &= 1{,}100{,}344. \notag
\end{align}

\step{5}{Total liabilities, long-term debt, equity}
\begin{align}
  TL &= (0.75)(8{,}000{,}000) \notag\\
  &= 6{,}000{,}000, \notag\\
  LTD &= TL - CL \notag\\
  &= (6{,}000{,}000) - (2{,}037{,}330) \notag\\
  &= 3{,}962{,}670, \notag\\
  E &= (8{,}000{,}000) - (6{,}000{,}000) \notag\\
  &= 2{,}000{,}000. \notag
\end{align}

\begin{center}
\begin{tabular}{lr@{\qquad}lr}
\toprule
Assets & \$ & Liabilities and equity & \$\\
\midrule
Cash & 1,100,000 & Accounts payable & 936,986\\
Accounts receivable & 2,297,059 & Short-term debt & 1,100,344\\
Inventory & 1,900,000 & Total current liabilities & 2,037,330\\
Total current assets & 5,297,059 & Long-term debt & 3,962,670\\
Net fixed assets & 2,702,941 & Shareholders' equity & 2,000,000\\
\midrule
Total assets & 8,000,000 & Total liabilities and equity & 8,000,000\\
\bottomrule
\end{tabular}
\end{center}

\intuition Each turnover ratio converts a flow (sales, COGS) into the stock it finances; once one
anchor (cash) pins down sales, the ratios fill in the balance sheet line by line, and the identity
$TA = TL + E$ closes it.

\bookans Matches the book's answer (pp.~403--404), line by line (figures rounded to the dollar).

\refbib Higgins Ch.~2, \S\,``Asset Turnover'' (collection period, inventory turnover, payables
period, days' sales in cash), pp.~42--48; ``Financial Leverage'' (current ratio, liabilities to
assets), pp.~48--53.

% =====================================================================
\exsec{15. Taco Mel: statements and cash flows for two years --- Ross 2.25--2.26, pp.~50--51}

\notation
\begin{itemize}
  \item Symbols as in Exercises 9 and 11; $\tau = 0.34$; $RE$: addition to retained earnings;
        $NP$: short-term notes (\pt{T\'itulos de d\'ivida de curto prazo}); $E$: owners' equity.
\end{itemize}
\assume Short-term notes are a current liability (inside NWC); \pt{Exig\'ivel a longo prazo} is
long-term debt; equity is the balancing item. Cash flow to creditors uses long-term debt only,
following Ross \S2.4.

\step{1}{Income statements (Problem 25)}
\begin{align}
  EBIT_{08} &= (7{,}233) - (2{,}487) - (591) - (1{,}038) \notag\\
  &= 3{,}117, \label{e15-ebit8}\\
  EBT_{08} &= (3{,}117) - (485) \notag\\
  &= 2{,}632, \notag\\
  Tax_{08} &= (0.34)(2{,}632) \notag\\
  &= 894.88, \notag\\
  NI_{08} &= (2{,}632) - (894.88) \notag\\
  &= 1{,}737.12, \notag\\
  RE_{08} &= (1{,}737.12) - (882) \notag\\
  &= 855.12, \notag\\
  EBIT_{09} &= (8{,}085) - (2{,}942) - (515) - (1{,}085) \notag\\
  &= 3{,}543, \label{e15-ebit9}\\
  EBT_{09} &= (3{,}543) - (579) \notag\\
  &= 2{,}964, \notag\\
  Tax_{09} &= (0.34)(2{,}964) \notag\\
  &= 1{,}007.76, \label{e15-tax9}\\
  NI_{09} &= (2{,}964) - (1{,}007.76) \notag\\
  &= 1{,}956.24, \notag\\
  RE_{09} &= (1{,}956.24) - (1{,}011) \notag\\
  &= 945.24. \label{e15-re9}
\end{align}

\step{2}{Balance sheets (Problem 25)}
\begin{align}
  CA_{08} &= (3{,}792) + (5{,}021) + (8{,}927) \notag\\
  &= 17{,}740, \notag\\
  TA_{08} &= (17{,}740) + (31{,}805) \notag\\
  &= 49{,}545, \notag\\
  CL_{08} &= (3{,}984) + (732) \notag\\
  &= 4{,}716, \notag\\
  E_{08} &= (49{,}545) - (4{,}716) - (12{,}700) \notag\\
  &= 32{,}129, \label{e15-e8}\\
  CA_{09} &= (4{,}041) + (5{,}892) + (9{,}555) \notag\\
  &= 19{,}488, \notag\\
  TA_{09} &= (19{,}488) + (33{,}921) \notag\\
  &= 53{,}409, \notag\\
  CL_{09} &= (4{,}025) + (717) \notag\\
  &= 4{,}742, \notag\\
  E_{09} &= (53{,}409) - (4{,}742) - (15{,}435) \notag\\
  &= 33{,}232. \label{e15-e9}
\end{align}

\begin{center}
\begin{tabular}{lrr@{\qquad}lrr}
\toprule
Income statement & 2008 & 2009 & Balance sheet & 2008 & 2009\\
\midrule
Sales & 7,233 & 8,085 & Cash & 3,792 & 4,041\\
COGS & 2,487 & 2,942 & Accounts receivable & 5,021 & 5,892\\
Other expenses & 591 & 515 & Inventory & 8,927 & 9,555\\
Depreciation & 1,038 & 1,085 & Current assets & 17,740 & 19,488\\
EBIT & 3,117 & 3,543 & Net fixed assets & 31,805 & 33,921\\
Interest & 485 & 579 & Total assets & 49,545 & 53,409\\
\cmidrule(lr){4-6}
Pretax income & 2,632 & 2,964 & Accounts payable & 3,984 & 4,025\\
Taxes (34\%) & 894.88 & 1,007.76 & Short-term notes & 732 & 717\\
Net income & 1,737.12 & 1,956.24 & Current liabilities & 4,716 & 4,742\\
Dividends & 882 & 1,011 & Long-term debt & 12,700 & 15,435\\
Add.\ to ret.\ earnings & 855.12 & 945.24 & Owners' equity & 32,129 & 33,232\\
\bottomrule
\end{tabular}
\end{center}

\step{3}{Cash flow from assets, 2009 (Problem 26)}
By \eqref{e9-ocfdef} with \eqref{e15-ebit9} and \eqref{e15-tax9},
\begin{align}
  OCF &= (3{,}543) + (1{,}085) - (1{,}007.76) \notag\\
  &= 3{,}620.24, \notag\\
  NCS &= (33{,}921 - 31{,}805) + (1{,}085) \notag\\
  &= 3{,}201, \notag\\
  \Delta NWC &= (19{,}488 - 4{,}742) - (17{,}740 - 4{,}716) \notag\\
  &= (14{,}746) - (13{,}024) \notag\\
  &= 1{,}722, \notag\\
  CFA &= (3{,}620.24) - (3{,}201) - (1{,}722) \notag\\
  &= -1{,}302.76. \label{e15-cfa}
\end{align}

\step{4}{Cash flow to creditors and shareholders, 2009 (Problem 26)}
\begin{align}
  CFC &= Int - (LTD_{09} - LTD_{08}) \notag\\
  &= (579) - (15{,}435 - 12{,}700) \notag\\
  &= (579) - (2{,}735) \notag\\
  &= -2{,}156, \label{e15-cfc}\\
  \text{net new equity} &= (E_{09} - E_{08}) - RE_{09} \notag\\
  &= (33{,}232 - 32{,}129) - (945.24) \notag\\
  &= 157.76, \notag\\
  CFS &= Div - \text{net new equity} \notag\\
  &= (1{,}011) - (157.76) \notag\\
  &= 853.24. \label{e15-cfs}
\end{align}
Check: $CFC + CFS = (-2{,}156) + (853.24) = -1{,}302.76 = CFA$ in \eqref{e15-cfa}.

\intuition Taco Mel's operations generated 3,620 but it spent 4,923 on fixed assets and working
capital, so investors as a group financed the gap of 1,303: creditors supplied 2,156 net (2,735 of
new long-term debt less 579 of interest), which covered both the gap and a net payout of 853 to
shareholders.

\refbib Ross Cap.~2, \S2.1 ``O balan\c{c}o patrimonial'', pp.~25--30; \S2.2, pp.~30--34;
\S2.4 ``Fluxo de caixa'', pp.~37--44.

\end{document}
